% TOP_PROBLEM: {"problemId": "problem.bartnik-cosmological-splitting-conjecture", "canonicalTitle": "Bartnik's cosmological splitting conjecture", "releaseRank": 398, "primaryDomain": "mathematical_physics", "primaryDomainLabel": "Mathematical physics", "displayStatus": "open_with_solved_subcases", "releaseStatus": "open_with_solved_subcases"}
% !TeX program = lualatex
% ATTEMPT_STATUS: unresolved
\documentclass[11pt]{article}
\usepackage[margin=25mm]{geometry}
\usepackage{fontspec}
\setmainfont{Latin Modern Roman}
\usepackage{amsmath,amssymb,mathtools}
\usepackage{unicode-math}
\setmathfont{Latin Modern Math}
\usepackage{xurl,hyperref}
\hypersetup{colorlinks=true,urlcolor=blue}
\setcounter{secnumdepth}{0}
\setlength{\parindent}{0pt}
\setlength{\parskip}{5pt}
\setlength{\emergencystretch}{3em}
\begin{document}
\section*{\texorpdfstring{Bartnik's cosmological splitting conjecture}{Bartnik's cosmological splitting conjecture}}\label{398}
\subsection{Definitions and mathematical statement}
Use Lorentzian signature $(-,+,\ldots,+)$. A Cauchy hypersurface meets every inextendible timelike curve exactly once, and global hyperbolicity supplies a well-behaved causal evolution. Let $(M,g)$ be a connected smooth time-oriented boundaryless spacetime of dimension $n\ge2$, globally hyperbolic with a compact Cauchy hypersurface, and complete along timelike geodesics in both time directions. Assume timelike convergence:
\[
 \operatorname{Ric}_g(v,v)\ge0\qquad\text{for every timelike }v.
\]
Bartnik's conjecture asks for a global isometry $(M,g)\cong({\mathbb{R}}\times S,-{\,\mathrm{d}} t^2+h)$ with $(S,h)$ compact Riemannian. The splitting is metric and global, not just the topological product supplied by global hyperbolicity. No timelike line or stationary symmetry is added as an assumption; obtaining the splitting without those extra hypotheses is the target.
\par\smallskip\textit{Formulation and concept references:} \href{https://arxiv.org/html/1810.03183v2}{[S1]}.

\subsection{Short English statement}
Must every timelike-complete spatially compact spacetime satisfying the timelike Ricci condition be a static metric product of time with a compact space?
\subsection{Research attempt: global Gaussian rigidity and the missing timelike line}
\textbf{Status.} The unrestricted splitting assertion is \textbf{UNRESOLVED} by this attempt. We prove it for metrics admitting a global proper-time Gaussian product, for Robertson--Walker metrics complete in the time coordinate, and for static metrics with compact spatial slice. These are additional geometric hypotheses, not consequences established from the catalogue assumptions.

\textbf{Current-source audit.} Section 1 of [S1] distinguishes the Lorentzian splitting theorem with a timelike line from Bartnik's conjecture without that line; its conformal-symmetry results impose further hypotheses. A June 2026 manuscript by McCann--Ohanyan entitled \emph{Positive resolution of Bartnik's cosmological splitting conjecture} is explicitly withdrawn in version 3, dated 22 June 2026 [E1]. The author notice identifies a false starting proposition and a counterexample of Ehrlich--Galloway. The still-visible original abstract is therefore not evidence that the conjecture has been settled. We use neither the withdrawn proposition nor a presumed global eikonal solution. This audit is important because the unverified global foliation needed below is exactly the sort of additional structure that can make a rigidity calculation deceptively easy.

Write $m=n-1\ge1$ for the spatial dimension. The convention for the Ricci tensor is the one in which a unit round Riemannian sphere has positive Ricci curvature and de Sitter spacetime has $\operatorname{Ric}=m g$.

\subsubsection{1. The product models and the full timelike condition}
For a product $g=-dt^2+h$ on $\mathbb R\times S$, the Levi-Civita connection preserves both factors and has no mixed curvature. Thus
\[
 \operatorname{Ric}_g(a\partial_t+w,a\partial_t+w)
 =\operatorname{Ric}_h(w,w).                            \tag{1}
\]
Every spatial vector $w$ can occur in a timelike vector by choosing $a^2>h(w,w)$. It follows that timelike convergence for the product is equivalent to $\operatorname{Ric}_h\ge0$, not merely to the vanishing of $\operatorname{Ric}_g(\partial_t,\partial_t)$.

If $S$ is compact, its Riemannian geodesics are complete. Product geodesics have affine time coordinate and a Riemannian geodesic as spatial component, so all timelike geodesics are complete in both directions. Each slice is Cauchy: along a future timelike curve, $t$ increases strictly and the spatial speed when parameterized by $t$ is less than one. If its time coordinate had a finite upper endpoint, compactness and this Lipschitz bound would give a limiting spatial point, allowing a timelike extension. The same holds in the past. Hence every inextendible timelike curve crosses every slice exactly once.

The vertical lines are globally maximizing. For a future timelike curve from $(t_0,x)$ to $(t_1,x)$,
\[
 L_g(\gamma)=\int_{t_0}^{t_1}\sqrt{1-|x'(t)|_h^2}\,dt
 \le t_1-t_0,
\]
with equality on the vertical segment. Thus a static product supplies a timelike line as well as a topological splitting. The converse direction in the conjecture must manufacture the analogous global maximizing object or an equivalent geometric structure.

\subsubsection{2. A scalar differential inequality with two-sided rigidity}
We first isolate the elementary analytic fact used repeatedly below.

\textbf{Lemma.} If a continuously differentiable function $H:\mathbb R\to\mathbb R$ satisfies
\[
 H'(t)\le-\frac1m H(t)^2,                              \tag{2}
\]
then $H\equiv0$.

Indeed $H$ is nonincreasing. If $H(t_0)<0$, it stays negative for $t\ge t_0$. On that half-line,
\[
 \left(\frac1H\right)'=-\frac{H'}{H^2}\ge\frac1m.
\]
Integrating to $t=t_0-m/H(t_0)$ forces $1/H(t)\ge0$, contrary to negativity and finiteness of $H(t)$. If $H(t_0)>0$, apply the preceding argument to $\widetilde H(s)=-H(-s)$, which obeys the same inequality. Both signs are impossible, proving the lemma.

This is a two-sided result. On a half-line, $H(t)=m/(t+c)$ solves the equality and need not vanish. It would be incorrect to infer rigidity from future existence alone. Equally, boundedness of a time coordinate is irrelevant unless that coordinate is affine proper time on the curves used in (2).

\subsubsection{3. A complete proof under a global Gaussian-coordinate hypothesis}
Assume, in addition to the conjecture's hypotheses, that the metric is globally
\[
 M=\mathbb R\times S,\qquad g=-dt^2+h_t,                \tag{3}
\]
where $h_t$ is a smooth family of Riemannian metrics and the coordinates cover all of $M$. This is much stronger than a smooth topological product. Define the shape endomorphism and its trace by
\[
 B^i{}_j=\tfrac12 h^{ik}\partial_t h_{kj},\qquad
 H=\operatorname{tr}B=\partial_t\log\sqrt{\det h_t}.
\]
A coordinate calculation gives
\[
 \Gamma^0_{ij}=\tfrac12\partial_t h_{ij},\qquad
 \Gamma^i_{0j}=B^i{}_j,\qquad \Gamma^\alpha_{00}=0.
\]
Substitution into the Ricci formula yields
\[
 \operatorname{Ric}_g(\partial_t,\partial_t)
 =-\partial_t H-\operatorname{tr}(B^2).                 \tag{4}
\]
Here $B$ is self-adjoint for $h_t$, so its eigenvalues are real and
$\operatorname{tr}(B^2)\ge H^2/m$ by Cauchy--Schwarz. Timelike convergence and (4) imply (2) for the function $t\mapsto H(t,x)$ at every fixed $x\in S$. The lemma forces $H(t,x)=0$ for every $(t,x)$.

Putting this back into (4),
\[
 0\le\operatorname{Ric}_g(\partial_t,\partial_t)
   =-\operatorname{tr}(B^2)\le0.
\]
Thus every eigenvalue of $B$ vanishes, so $\partial_t h_t=0$. The metric in (3) is the required product, and (1) also shows $\operatorname{Ric}_h\ge0$.

This proof needs neither a spatial integration nor a pointwise assumption on the sign of the initial mean curvature. Completeness of the coordinate normal geodesics is encoded by the fact that (3) is valid for all $t\in\mathbb R$. Compactness is needed for the catalogue's compact factor but not for this local-in-space Riccati argument.

What is not proved is (3). Normal geodesics launched from a Cauchy hypersurface can encounter focal points, and even before a focal point a normal exponential map can fail to be globally injective. Timelike geodesic completeness says the individual geodesics continue; it does not say their normal coordinate chart continues as a one-to-one smooth foliation. At a caustic, $B$ and $H$ for that chart need not remain smooth, so the lemma cannot be applied across it by assertion.

\subsubsection{4. Warped products make the completeness issue transparent}
Consider
\[
 g=-dt^2+a(t)^2h,\qquad a(t)>0,
\]
on $I\times S$, with $I$ an open interval. Vertical curves are unit-speed timelike geodesics, so timelike completeness forces $I=\mathbb R$ in this particular coordinate model. Equation (4), or a direct computation, gives
\[
 \operatorname{Ric}_{tt}=-m\frac{a''}{a}.               \tag{5}
\]
Timelike convergence implies $a''\le0$. A positive concave function on the entire real line is constant. To prove this, if $a'(t_0)>0$, concavity gives $a'(t)\ge a'(t_0)$ for $t<t_0$, hence
$a(t)\le a(t_0)-a'(t_0)(t_0-t)$, negative for sufficiently negative $t$. If $a'(t_0)<0$, the same argument in the future gives negativity. Thus every derivative is zero.

The conclusion is again a static product. But $a(t)=t$ on $(0,\infty)$ is positive, concave and nonconstant; its vertical timelike geodesics reach $t=0$ in finite proper time. This stress test shows exactly which completeness direction the elementary proof uses.

For $a(t)=\cosh t$, (5) is $-m<0$. The de Sitter warped product is therefore not a counterexample under timelike convergence, despite its compact spatial slices and timelike completeness. Testing only scalar curvature, or imposing the opposite sign convention for timelike Ricci without changing the hypothesis, would reverse the conclusion incorrectly.

More generally, if a model is written
$-N(t)^2dt^2+a(t)^2h$ with positive lapse, one must first use proper time
$s=\int N(t)\,dt$. Vertical timelike completeness requires that the range of $s$ be all of $\mathbb R$. Having a coordinate $t\in\mathbb R$ is not enough when that integral has a finite endpoint.

\subsubsection{5. A static lapse on a compact slice is forced to be constant}
There is a different special case in which global Gaussian coordinates are not initially assumed. Let
\[
 g=-N(x)^2dt^2+h(x),\qquad N>0,
\]
on $\mathbb R\times S$, with $S$ closed and connected. The nonzero connection coefficients involving time are
\[
 \Gamma^0_{0i}=\partial_i\log N,\qquad
 \Gamma^i_{00}=N h^{ij}\partial_jN.
\]
A direct Ricci computation gives
\[
 \operatorname{Ric}_{00}=N\Delta_hN,\qquad
 \operatorname{Ric}_{ij}=\operatorname{Ric}_{h,ij}
                         -N^{-1}\operatorname{Hess}_{ij}N,
 \qquad \operatorname{Ric}_{0i}=0.                     \tag{6}
\]
Use the Laplacian convention $\Delta=\operatorname{div}\nabla$. Since $\partial_t$ is timelike, convergence implies $\Delta_hN\ge0$. Integration over the closed slice gives $\int_S\Delta_hN=0$, hence $\Delta_hN=0$ everywhere. Integration by parts then gives
$\int_S|\nabla N|^2=-\int_SN\Delta_hN=0$. Connectedness makes $N$ constant, and rescaling $t$ produces $-ds^2+h$.

Compactness without boundary is essential in this proof: neither the vanishing integral nor the absence of boundary terms is automatic on a noncompact or bounded slice. Nor is an arbitrary stationary metric static. A general stationary metric has a shift one-form; (6) was computed only with that shift zero. Therefore this special case does not supply the timelike Killing field or its hypersurface orthogonality in the unrestricted conjecture.

\subsubsection{6. A parallel field gives a product under an exactness hypothesis}
Another precise sufficient condition is useful for locating the global issue. Suppose $V$ is a globally defined parallel unit timelike vector field and the one-form
$\alpha=-g(V,\cdot)$ is exact, say $\alpha=dt$. Then $dt(V)=1$, and timelike geodesic completeness makes the integral curves of $V$ complete: parallelness gives $\nabla_VV=0$ and constant speed.

The flow maps any nonempty level $S=t^{-1}(0)$ diffeomorphically to the levels $t^{-1}(s)$. Every point reaches $S$ after flowing for time $-t$, so
\[
 \Psi:\mathbb R\times S\longrightarrow M,\qquad
 \Psi(s,x)=\operatorname{Fl}_V^s(x)
\]
is a diffeomorphism, with explicit inverse obtained by the same flow. Parallelness implies $\mathcal L_Vg=0$, while $TS=\ker\alpha$ is orthogonal to $V$. Hence $\Psi^*g=-ds^2+h$.

If there is a compact Cauchy hypersurface $C$, the projection along the $V$-orbits maps $C$ bijectively onto $S$: every complete timelike orbit meets $C$ exactly once. Locally this projection is a diffeomorphism by transversality. Thus $S$ is diffeomorphic to $C$ and is compact. This proves the desired global compact product under the two extra assumptions just stated.

A parallel field only makes $\alpha$ closed. Exactness is a separate global assertion and has not been silently inferred from closedness. The calculation is a sufficient criterion, not a claim that the chosen field in every quotient example must admit this exact potential. In particular, replacing a Lorentzian splitting theorem by a local parallel-field argument requires checking global topology.

\subsubsection{7. The timelike-line route and its first gap}
The external Lorentzian splitting theorem used in the discussion of [S1] says, in the smooth timelike-complete setting with timelike convergence and a timelike line, that the spacetime splits as a Lorentzian metric product. A timelike line is an inextendible timelike geodesic each of whose segments maximizes Lorentzian distance between its endpoints. The catalogue deliberately does not assume one.

One may take increasingly long maximizing segments whose endpoints escape to the past and future, and use compactness of a Cauchy slice to recenter them. There are two distinct compactness questions. Their crossing points have a convergent subsequence on the compact slice. Their unit future timelike tangent vectors need not have a convergent subsequence in the unit timelike bundle: the unit hyperboloid is noncompact.

In flat coordinates, the future unit vectors
\[
 v_j=\cosh(j)\,\partial_t+\sinh(j)\,\partial_{x_1}
\]
have $g(v_j,v_j)=-1$ but escape every compact subset of the tangent space. Their projective directions tend to a null direction. Thus a limit-curve argument without control of this boost parameter may yield a null line rather than a timelike one. This elementary example concerns the compactness inference, not a counterexample spacetime.

The source's no-observer-horizon hypothesis is designed to exclude such null-line behavior and then invoke splitting. It cannot be inserted into the proof as though it followed from compact Cauchy surfaces alone. Likewise the existence of a complete timelike conformal Killing field is a substantive extra assumption in [S1].

\subsubsection{8. Verification and conclusion}
The checks are deliberately structural: (4) specializes to (5); the trace inequality holds also for spatial dimension one; a constant lapse reduces (6) to (1); and the examples $a=t$ and $a=\cosh t$ violate respectively completeness and timelike convergence. All time parameters in the rigidity arguments are proper time or are explicitly converted to it. No inference from causal completeness to a globally injective normal exponential map has been made.

The strongest results proved here are complete splitting under a global all-time Gaussian form or a static compact-slice form. The first unresolved step for the original hypotheses is the construction of a timelike line, or another global structure strong enough to justify an all-time smooth Riccati argument. The withdrawn 2026 announcement does not supply that step. The catalogue target remains \textbf{UNRESOLVED} by this attempt.

\subsection{Sources}
\begin{itemize}
\item[S1]Costa e Silva–Flores–Herrera: Conformal symmetries and Bartnik splitting.
\url{https://arxiv.org/html/1810.03183v2}.
\textit{Formulation source.}
\item[E1]R. J. McCann and A. Ohanyan, \emph{Positive resolution of Bartnik's cosmological splitting conjecture}, arXiv:2606.03873v3, 22 June 2026, withdrawn.
\url{https://arxiv.org/abs/2606.03873v3}.
\textit{Withdrawal notice and author explanation inspected 27 September 2026. Cited only for current-status hygiene, not as a splitting theorem.}
\end{itemize}
\end{document}
